% GENERATED FILE. Edit canonical lessons, then run npm run generate:books.
% canonical-source-hash: fnv1a64:9030707c3f3a2d5f
% canonical-lessons: KA-C505-sraddhe, KA-C505-nirlaksya, KA-C505-cata, KA-C505-sramisu, KA-C505-bedu

\chapter{Effort and Neglect}
\label{ch:ka-effort-and-neglect}
\hlchaptermodality{505}

\begin{chapteropening}
\textbf{By the end of this chapter:} \emph{I can say dedication, neglect, an addiction, to strive and to beg.}

Complete the last lesson of chapter 505: I can say dedication, neglect, an addiction, to strive and to beg.
\end{chapteropening}

\section[\texorpdfstring{\kn{ಶ್ರದ್ಧೆ} (śraddhe)}{ಶ್ರದ್ಧೆ (śraddhe)}]{\kn{ಶ್ರದ್ಧೆ} (śraddhe) — dedication; devotion}

\label{lesson:KA-C505-sraddhe}



\begin{quote}
Before the new one: say the Kannada for a coward, then the Kannada for a stammer.
\end{quote}

\subsection*{You'll want to know: \kn{ಶ್ರದ್ಧೆ}}
\textbf{\kn{ಶ್ರದ್ಧೆ}} — \emph{śraddhe} — \textquotedblleft{}dedication; devotion\textquotedblright{}.

\begin{quote}
\kn{ಅವಳು} \kn{ಶ್ರದ್ಧೆಯಿಂದ} \kn{ಸಂಗೀತ} \kn{ಕಲಿಯುತ್ತಾಳೆ}.
\end{quote}

\begin{quote}
\emph{avaḷu śraddheyinda saṅgīta kaliyuttāḷe.}
\end{quote}

\begin{quote}
She learns music with dedication.
\end{quote}

\kn{ಶ್ರದ್ಧೆಯಿಂದ}, wholeheartedly.

\begin{cousinweb}
\textbf{\kn{ಶ್ರದ್ಧೆ}} (\emph{śraddhe}) is borrowed from Sanskrit \emph{śraddhā}, \textquotedblleft{}trust, faith\textquotedblright{} (CDIAL 12678); in Kannada it is steady devotion to a task.
\end{cousinweb}

\begin{grammarlens}[title={What you've built}]
One more word for describing what you lived and felt.
\end{grammarlens}

\subsection*{Your turn}
\begin{itemize}
\raggedright
  \item \emph{Say it:} \emph{śraddhe}
  \item \emph{Say it:} \emph{śraddhe}, once more
  \item \emph{Say it:} \emph{avaḷu śraddheyinda saṅgīta kaliyuttāḷe.} — She learns music with dedication.
  \item \emph{Recall:} say the Kannada for a coward, then the Kannada for a stammer, then say \emph{śraddhe} again
\end{itemize}

\begin{tcolorbox}[breakable,colback=teal!4,colframe=teal!35!black,title={Before you move on}]
What does \kn{ಶ್ರದ್ಧೆ} mean? (\textquotedblleft{}Dedication; devotion\textquotedblright{}.) Say it once more.
\end{tcolorbox}
\section[\texorpdfstring{\kn{ನಿರ್ಲಕ್ಷ್ಯ} (nirlakṣya)}{ನಿರ್ಲಕ್ಷ್ಯ (nirlakṣya)}]{\kn{ನಿರ್ಲಕ್ಷ್ಯ} (nirlakṣya) — neglect; carelessness}

\label{lesson:KA-C505-nirlaksya}



\begin{quote}
Before the new one: say the Kannada for a stammer, then the Kannada for dedication.
\end{quote}

\subsection*{You'll want to know: \kn{ನಿರ್ಲಕ್ಷ್ಯ}}
\textbf{\kn{ನಿರ್ಲಕ್ಷ್ಯ}} — \emph{nirlakṣya} — \textquotedblleft{}neglect; carelessness\textquotedblright{}.

\begin{quote}
\kn{ಚಾಲಕನ} \kn{ನಿರ್ಲಕ್ಷ್ಯದಿಂದ} \kn{ಅಪಘಾತ} \kn{ಆಯಿತು}.
\end{quote}

\begin{quote}
\emph{cālakana nirlakṣyadinda apaghāta āyitu.}
\end{quote}

\begin{quote}
The accident happened through the driver's carelessness.
\end{quote}

The verb is \kn{ನಿರ್ಲಕ್ಷಿಸು}, to ignore.

\begin{cousinweb}
\textbf{\kn{ನಿರ್ಲಕ್ಷ್ಯ}} (\emph{nirlakṣya}) is built from Sanskrit \emph{nir-}, \textquotedblleft{}without\textquotedblright{}, and \emph{lakṣya}, \textquotedblleft{}what is to be marked or watched\textquotedblright{}; Monier-Williams lists \emph{nirlakṣya} as \textquotedblleft{}unobservable\textquotedblright{} (Monier-Williams). Kannada uses it for failing to watch.
\end{cousinweb}

\begin{grammarlens}[title={What you've built}]
One more word for describing what you lived and felt.
\end{grammarlens}

\subsection*{Your turn}
\begin{itemize}
\raggedright
  \item \emph{Say it:} \emph{nirlakṣya}
  \item \emph{Say it:} \emph{nirlakṣya}, once more
  \item \emph{Say it:} \emph{cālakana nirlakṣyadinda apaghāta āyitu.} — The accident happened through the driver's carelessness.
  \item \emph{Recall:} say the Kannada for a stammer, then the Kannada for dedication, then say \emph{nirlakṣya} again
\end{itemize}

\begin{tcolorbox}[breakable,colback=teal!4,colframe=teal!35!black,title={Before you move on}]
What does \kn{ನಿರ್ಲಕ್ಷ್ಯ} mean? (\textquotedblleft{}Neglect; carelessness\textquotedblright{}.) Say it once more.
\end{tcolorbox}
\section[\texorpdfstring{\kn{ಚಟ} (caṭa)}{ಚಟ (caṭa)}]{\kn{ಚಟ} (caṭa) — an addiction; a bad habit}

\label{lesson:KA-C505-cata}



\begin{quote}
Before the new one: say the Kannada for dedication, then the Kannada for neglect.
\end{quote}

\subsection*{You'll want to know: \kn{ಚಟ}}
\textbf{\kn{ಚಟ}} — \emph{caṭa} — \textquotedblleft{}an addiction; a bad habit\textquotedblright{}.

\begin{quote}
\kn{ಅವನಿಗೆ} \kn{ಸಿಗರೇಟಿನ} \kn{ಚಟ} \kn{ಇದೆ}.
\end{quote}

\begin{quote}
\emph{avanige sigarēṭina caṭa ide.}
\end{quote}

\begin{quote}
He has a smoking habit.
\end{quote}

\kn{ಚಟ} \kn{ಬಿಡು}, to kick a habit.

\begin{cousinweb}
\textbf{\kn{ಚಟ}} (\emph{caṭa}) is an everyday Kannada word for a habit that has a grip on you; its origin is not settled.
\end{cousinweb}

\begin{grammarlens}[title={What you've built}]
One more word for describing what you lived and felt.
\end{grammarlens}

\subsection*{Your turn}
\begin{itemize}
\raggedright
  \item \emph{Say it:} \emph{caṭa}
  \item \emph{Say it:} \emph{caṭa}, once more
  \item \emph{Say it:} \emph{avanige sigarēṭina caṭa ide.} — He has a smoking habit.
  \item \emph{Recall:} say the Kannada for dedication, then the Kannada for neglect, then say \emph{caṭa} again
\end{itemize}

\begin{tcolorbox}[breakable,colback=teal!4,colframe=teal!35!black,title={Before you move on}]
What does \kn{ಚಟ} mean? (\textquotedblleft{}An addiction; a bad habit\textquotedblright{}.) Say it once more.
\end{tcolorbox}
\section[\texorpdfstring{\kn{ಶ್ರಮಿಸು} (śramisu)}{ಶ್ರಮಿಸು (śramisu)}]{\kn{ಶ್ರಮಿಸು} (śramisu) — to strive; to work hard}

\label{lesson:KA-C505-sramisu}



\begin{quote}
Before the new one: say the Kannada for neglect, then the Kannada for an addiction.
\end{quote}

\subsection*{You'll want to know: \kn{ಶ್ರಮಿಸು}}
\textbf{\kn{ಶ್ರಮಿಸು}} — \emph{śramisu} — \textquotedblleft{}to strive; to work hard\textquotedblright{}.

\begin{quote}
\kn{ರೈತರು} \kn{ಬೆಳೆಗಾಗಿ} \kn{ಹಗಲು} \kn{ರಾತ್ರಿ} \kn{ಶ್ರಮಿಸುತ್ತಾರೆ}.
\end{quote}

\begin{quote}
\emph{raitaru beḷegāgi hagalu rātri śramisuttāre.}
\end{quote}

\begin{quote}
Farmers toil day and night for the crop.
\end{quote}

Past (I …): \kn{ಶ್ರಮಿಸಿದೆ}, \emph{śramiside}.

For what: -\kn{ಗಾಗಿ}. \kn{ಶ್ರಮ} is toil.

\begin{cousinweb}
\textbf{\kn{ಶ್ರಮಿಸು}} (\emph{śramisu}) is built on Sanskrit \emph{śrama}, \textquotedblleft{}labour, fatigue\textquotedblright{} (CDIAL 12682), with the verb ending -\emph{isu}.
\end{cousinweb}

\begin{grammarlens}[title={What you've built}]
One more word for describing what you lived and felt.
\end{grammarlens}

\subsection*{Your turn}
\begin{itemize}
\raggedright
  \item \emph{Say it:} \emph{śramisu}
  \item \emph{Say it:} \emph{śramisu}, once more
  \item \emph{Say it:} \emph{raitaru beḷegāgi hagalu rātri śramisuttāre.} — Farmers toil day and night for the crop.
  \item \emph{Recall:} say the Kannada for neglect, then the Kannada for an addiction, then say \emph{śramisu} again
\end{itemize}

\begin{tcolorbox}[breakable,colback=teal!4,colframe=teal!35!black,title={Before you move on}]
What does \kn{ಶ್ರಮಿಸು} mean? (\textquotedblleft{}To strive; to work hard\textquotedblright{}.) Say it once more.
\end{tcolorbox}
\section[\texorpdfstring{\kn{ಬೇಡು} (bēḍu)}{ಬೇಡು (bēḍu)}]{\kn{ಬೇಡು} (bēḍu) — to beg; to plead}

\label{lesson:KA-C505-bedu}



\begin{quote}
Before the new one: say the Kannada for an addiction, then the Kannada for to strive.
\end{quote}

\subsection*{You'll want to know: \kn{ಬೇಡು}}
\textbf{\kn{ಬೇಡು}} — \emph{bēḍu} — \textquotedblleft{}to beg; to plead\textquotedblright{}.

\begin{quote}
\kn{ಮುದುಕ} \kn{ದಾರಿಯಲ್ಲಿ} \kn{ಊಟಕ್ಕಾಗಿ} \kn{ಬೇಡುತ್ತಿದ್ದನು}.
\end{quote}

\begin{quote}
\emph{muduka dāriyalli ūṭakkāgi bēḍuttiddanu.}
\end{quote}

\begin{quote}
The old man was begging for food by the road.
\end{quote}

Past (I …): \kn{ಬೇಡಿದೆ}, \emph{bēḍide}.

What you beg for takes -\kn{ಗಾಗಿ} or -\kn{ಅನ್ನು}. \kn{ಭಿಕ್ಷೆ} \kn{ಬೇಡು}, to beg for alms.

\begin{cousinweb}
\textbf{\kn{ಬೇಡು}} (\emph{bēḍu}) is native Dravidian: DEDR 5528 lists Kannada \emph{bēḍu}, \textquotedblleft{}to wish, desire, beg, entreat, request\textquotedblright{}, in the same article as \kn{ಬೇಕು} and \kn{ಬೇಡ}.
\end{cousinweb}

\begin{grammarlens}[title={What you've built}]
One more word for describing what you lived and felt.
\end{grammarlens}

\subsection*{Your turn}
\begin{itemize}
\raggedright
  \item \emph{Say it:} \emph{bēḍu}
  \item \emph{Say it:} \emph{bēḍu}, once more
  \item \emph{Say it:} \emph{muduka dāriyalli ūṭakkāgi bēḍuttiddanu.} — The old man was begging for food by the road.
  \item \emph{Recall:} say the Kannada for an addiction, then the Kannada for to strive, then say \emph{bēḍu} again
\end{itemize}

\begin{tcolorbox}[breakable,colback=teal!4,colframe=teal!35!black,title={Before you move on}]
What does \kn{ಬೇಡು} mean? (\textquotedblleft{}To beg; to plead\textquotedblright{}.) Say it once more.
\end{tcolorbox}
